(1) 标准动能为 \(T=I(\dot\theta^2+\dot\phi^2\sin^2\theta)/2+I_3(\dot\psi+\dot\phi\cos\theta)^2/2\)，重力势 \(V=mg\ell\cos\theta\)。消去循环自旋坐标后，\(\theta\) 的 Euler--Lagrange 方程在 \(\dot\theta=0,\dot\phi=\Omega,\omega_3=\mathrm{const}\) 时化为 \(0=\sin\theta[I\Omega^2\cos\theta-I_3\omega_3\Omega+mg\ell]\)。因倾角非零，括号必须为零。

(2) 这是关于 \(\Omega\) 的二次方程，根为 \(\Omega_\pm=\{I_3\omega_3\pm[(I_3\omega_3)^2-4Img\ell\cos\theta]^{1/2}\}/(2I\cos\theta)\)。减号为高速自旋下的慢进动，正号为快进动；两根的和 \(I_3\omega_3/(I\cos\theta)\)，积 \(mg\ell/(I\cos\theta)\)，量纲均为角频率的适当组合。自旋速率 \(\omega_3\) 不应误当作欧拉角 \(\dot\psi\)。

(3) 判别式非负给 \((I_3\omega_3)^2\ge4Img\ell\cos\theta\)，即 \(\omega_3\ge 2\sqrt{Img\ell\cos\theta}/I_3\)。等号对应快慢匀进动根并合，低于此阈值不能以所指定的 \(\theta\)、方向和纯匀角速度实现解。此判据是匀进动存在性条件，并不自动保证扰动稳定。

(4) 对 \(I_3^2\omega_3^2\gg4Img\ell\cos\theta\) 使用 \(\sqrt{1-\epsilon}=1-\epsilon/2+O(\epsilon^2)\)，得 \(\Omega_-\simeq mg\ell/(I_3\omega_3)\)、\(\Omega_+\simeq I_3\omega_3/(I\cos\theta)-mg\ell/(I_3\omega_3)\)。慢进动符合角动量向量以 \(\boldsymbol\tau=\boldsymbol r\times m\boldsymbol g\) 所确定的方向旋转；如果颠倒势能的 \(mg\ell\cos\theta\) 符号，将产生错误的进动方向。
